% Generated draft: requires mathematical and source review.
\documentclass[11pt]{article}
\usepackage[margin=25mm]{geometry}
\usepackage{fontspec}
\setmainfont{Latin Modern Roman}
\usepackage{amsmath,amssymb,mathtools,unicode-math}
\setmathfont{Latin Modern Math}
\usepackage{enumitem,needspace,xurl,hyperref}
\setlength{\parindent}{0pt}
\setlength{\parskip}{5pt}
\setlength{\emergencystretch}{3em}
\newcommand{\N}{\mathbb{N}}
\newcommand{\Z}{\mathbb{Z}}
\newcommand{\Q}{\mathbb{Q}}
\newcommand{\R}{\mathbb{R}}
\newcommand{\C}{\mathbb{C}}
\newcommand{\F}{\mathbb{F}}
\newcommand{\A}{\mathbb{A}}
\newcommand{\PP}{\mathbb{P}}
\newcommand{\E}{\mathbb{E}}
\newcommand{\eps}{\varepsilon}
\newcommand{\dd}{\,\mathrm{d}}
\begin{document}
\section*{UnsolvedMath ID: 30003853}
\textbf{Source problem code:} catalogue entry
\subsection{Definitions and mathematical statement}
Let \(F\) denote Thompson's group in its standard interval realization: the group, under composition, of increasing piecewise linear homeomorphisms \(f:[0,1]\to[0,1]\) with finitely many breakpoints, all in \(\Z[1/2]\cap(0,1)\), and with every linear slope in \(\{2^k:k\in\Z\}\). Here \(\Z[1/2]=\{m/2^j:m\in\Z,\ j\in\Z_{\geq0}\}\).

A subgroup \(H\leq F\) is finitely presented if, as an abstract group, it admits a presentation with finitely many generators and finitely many defining relators. Its commutator subgroup is
\[
[H,H]=\langle xyx^{-1}y^{-1}:x,y\in H\rangle,
\]
and its abelianization is \(H_{\mathrm{ab}}=H/[H,H]\). An abelian group \(A\), written additively, is torsion-free if \(na=0\) implies \(a=0\) for every \(a\in A\) and every integer \(n\geq1\).

The question is whether every finitely presented subgroup \(H\leq F\) has torsion-free abelianization. Equivalently, is it true that, for every finitely presented \(H\leq F\), every \(h\in H\), and every integer \(n\geq1\),
\[
h^n\in[H,H]\quad\Longrightarrow\quad h\in[H,H]?
\]

The supplied record labels this question open. Its supplied literature assessment attributes to Rover an example disproving the analogous assertion for finitely generated subgroups, while stating that the report leaves the finitely presented case open. These are source qualifications, not independently verified conclusions. The report's bibliographic year is supplied as 2019, whereas its report designation and proposed year refer to 2018; this draft does not resolve that dating distinction.
\subsection{Short English statement}
Must the abelianization of every finitely presented subgroup of Thompson's group \(F\) contain no nonzero element of finite order? The supplied record distinguishes this question from the finitely generated version, which it reports is disproved by Rover's example. Its open-status assessment requires independent review.
\subsection{Sources}
\begin{itemize}
\item Supplied local UnsolvedMath registry record, permanent ID \(30003853\), source problem code \texttt{catalogue entry}, titled ``Abelianizations of Finitely Presented Thompson \(F\) Subgroups,'' in the Oberwolfach Reports Open Problems collection. The proposed year is supplied as 2018; no individual proposer is identified. The record labels the statement \texttt{corrected\_verified} and reports a manual semantic, typographic, and LaTeX review; that supplied label does not constitute mathematical review of this draft.
\item Supplied report citation: \emph{Cohomological and Metric Properties of Groups of Homeomorphisms of R} (2019), Oberwolfach Reports, pp. 1579--1633, DOI 10.4171/owr/2018/26. \url{https://doi.org/10.4171/owr/2018/26}. The record describes Oberwolfach Reports as documenting mathematical workshops and their open-problem discussions.
\item Supplied literature source: ``Oberwolfach Report 26/2018.'' \url{https://ems.press/content/serial-article-files/46748?nt=1}.
\item Literature assessment dated exactly \texttt{2026-08-21}: the supplied curation labels the question open, attributes the failure of the finitely generated analogue to Rover's example, states that the report explicitly leaves the finitely presented case open, and reports that no later primary proof or counterexample was found. This is supplied dated curation derived from the OWR material, not independently verified source consultation or a new literature assessment. Time-sensitive status and cited claims require independent verification.
\end{itemize}
\end{document}
